Le Strontium ${ }_{38}^{90} \mathrm{Sr}$ est l'un des principaux noyaux se trouvant dans les déchets des réacteurs nucléaires électrogènes. ${ }_{38}^{90} \mathrm{Sr}$ se désintègre en Yttrium ${ }_{Z}^{A} \boldsymbol{Y}$ en émettant une particule $\boldsymbol{\beta}^{-}$. Il est très nocif quand la particule est inhalée ou ingérée. Le noyau ${ }_{38}^{88} S r$ est l'un des isotopes de l'élément Strontium.

\section*{Données:}
\begin{center}
\begin{tabular}{|l|l|l|l|l|l|}
\hline
Particule & ${ }_{38}^{90} \mathrm{Sr}$ & ${ }_{Z}^{A} Y$ & Électron & Proton & Neutron \\
\hline
Masse en ( $u$ ) & ,90773 & 89,90714 & 0,0005 & 1,00728 & 1,00866 \\
\hline
\multicolumn{4}{|l|}{Demi-vie de ${ }_{38}^{90} \mathrm{Sr}: \quad t_{1 / 2}=28,9$ ans} & \multicolumn{2}{|c|}{$1 u=931,5 \mathrm{MeV} . c^{-2}$} \\
\hline
\multicolumn{2}{|l|}{Energie de liaison du noyau ${ }_{38}^{88} \mathrm{Sr}$ :} & \multicolumn{2}{|r|}{$E_{\ell}\left({ }_{38}^{88} S r\right)=768,47 \mathrm{MeV}$} & \multicolumn{2}{|c|}{1 an $=365$ jours} \\
\hline
\end{tabular}
\end{center}

\begin{enumerate}
  \item Écrire l'équation de désintégration de ${ }_{38}^{90} S r$ en précisant les valeurs de $A$ et $Z$ du noyau fils ${ }_{Z}^{A} Y$.
  \item Calculer, en unité MeV , l'énergie de liaison $E_{\ell}\left({ }_{38}^{90} \mathrm{Sr}\right)$ du noyau ${ }_{38}^{90} \mathrm{Sr}$.
  \item Déterminer, en justifiant, le noyau le plus stable parmi ${ }_{38}^{90} S r$ et ${ }_{38}^{88} S r$.
  \item Un échantillon de Strontium 90 a une activité initiale $a_{0}=5,11.10^{12} \mathrm{~Bq}$ à l'instant $t_{0}=0$.\\
4.1. Calculer le nombre $N_{0}$ des noyaux présents initialement dans cet échantillon.\\
4.2. Déterminer l'instant $t_{1}$ pour lequel l'activité de cet échantillon devient $a_{1}=10 \% a_{0}$.
\end{enumerate}


